%%%%%%%%%%  scan idx 76  |  folio 63  %%%%%%%%%%
$1$
pour $F=\cO_{(x)}$ pour tout $x\in X$.

Ainsi, $A(X)$ et $G(\underline K(X))$ ont dans le cas envisagé même rang fini en tout degré, et $G(\underline K(X))$ est sans torsion, puisqu'il en est ainsi de $A(X)$. De l'assertion b) ci-dessus résulte, de manière «~formelle~», que l'application $\Ctil$ de $\underline K(X)$ dans $\widetilde{A(X)}$ est surjective, modulo des groupes finis. Il en résulte facilement, puisque $G(\underline K(X))$ et $A(X)$ ont même rang fini en tout degré, que l'homomorphisme $G(\Ctil)$ de $G(\underline K(X))$ dans $G(\widetilde{A(X)})$ est bijectif modulo des groupes finis en tout degré (on utilisera une récurrence sur le degré). Comme $G(\underline K(X))$ est sans torsion, cet homomorphisme $G(\Ctil)$ est donc injectif.
\begin{flushright}cqfd.\end{flushright}

N.B. L'homomorphisme $\Ctil$ n'est pas surjectif en général, même si $X$ est un espace projectif de dimension $2$.

\subsection*{\S 5. Définition faisceautique des classes de Chern. Application à l'étude des morphismes d'injection.}

\textbf{Théorème 2.16.} \emph{Soit $X$ une variété quasi-projective non singulière. Alors, pour tout $x\in\underline K(X)$, et pour tout entier $i\geq0$, l'élément $\gamma^i(x-\varepsilon(x))$ est de filtration $\geq i$. Soit $c^i(x)\in G^i(\underline K(X))$ la classe de cet élément modulo $\underline K(X)^{i+1}$~; si l'on pose~:}
\[
 \widetilde c(x)=\left[\varepsilon(x),1+\sum_{i\geq0}c^i(x)\right]\in\widetilde{G(\underline K(X))},
\]
\emph{alors $\widetilde c$ satisfait aux conditions analogues à celles du théorème 2.2.}

La première assertion résulte du théorème 2.14, ainsi que le fait que $x\longrightarrow\widetilde c(x)$ soit un homomorphisme de $\lambda$-anneaux (la formule (1.25) implique $\gamma_t(x+y)=\gamma_t(x)\gamma_t(y)$, ce qui montre déjà que $\widetilde c$ est un homomorphisme additif). La «~fonctorialité~» de $\widetilde c$ se vérifie immédiatement, et 
